\documentclass[11pt,a4paper]{article}

\usepackage[T1]{fontenc}
\usepackage[utf8]{inputenc}
\usepackage[italian,english]{babel}
\usepackage{amsmath,amssymb}
\usepackage{booktabs}
\usepackage{geometry}
\geometry{margin=2.6cm}
\usepackage{xcolor}
\usepackage[colorlinks=true,linkcolor=blue,urlcolor=blue]{hyperref}
\usepackage{listings}
\usepackage{enumitem}

\lstset{
  basicstyle=\ttfamily\footnotesize,
  breaklines=true,
  frame=single,
  framesep=4pt,
  rulecolor=\color{gray!50},
  backgroundcolor=\color{gray!5},
  columns=fullflexible,
}

\newcommand{\MRE}{\mathrm{MRE}}
\newcommand{\TODO}[1]{\textcolor{red}{\textbf{[TODO: #1]}}}
\newcommand{\OPEN}[1]{\textcolor{orange!80!black}{\textbf{[APERTO: #1]}}}

\title{\textbf{Il pavimento di rumore dell'MRE in GibbsPCD:\\
diagnosi di un errore in arXiv:2603.27312}}
\author{Nota tecnica interna --- GSP / UrbIA}
\date{27 luglio 2026 \quad --- \quad versione 0.1}

\begin{document}
\selectlanguage{italian}
\maketitle

\begin{abstract}
\noindent
Il paper GibbsPCD (arXiv:2603.27312, sottomesso ad ACM TKDD) enuncia il
pavimento di rumore dell'errore relativo medio come $\MRE_{\min}\approx
1/(2\sqrt{N})$. Questa espressione è il \emph{massimo errore assoluto}
di una proporzione stimata su $N$ estrazioni, non un pavimento
dell'errore \emph{relativo}: essendo gli $\alpha_j$ dei vincoli
censuari tipicamente $\ll 1/2$, il pavimento corretto risulta una o due
decadi più alto. Si documenta la derivazione corretta, comprensiva del
fattore $\sqrt{2/\pi}$ di conversione fra deviazione standard e scarto
assoluto medio, e la sua validazione empirica su due popolazioni
sintetiche indipendenti (accordo entro l'1--2\%). Si delimita l'impatto
sul paper: le affermazioni sul \emph{tasso} $O(1/\sqrt{N})$ restano
valide, cade invece il potere discriminante del test ``$\MRE$ osservato
sopra il pavimento $\Rightarrow$ domina il mixing bias''.
Lo stato è \textbf{diagnosi completata, correzione non ancora decisa}:
si attende il referee report.
\end{abstract}

\tableofcontents

% ======================================================================
\section{Stato e contesto}

\paragraph{Origine.} Sessione di lavoro del 27 luglio 2026 sulla
generalizzazione della pipeline GSP a nuovi comuni (Parma, 034027). Il
fit K7C su Parma ha prodotto una coppia di diagnostiche
(esatto + Gibbs sullo stesso constraint set) che ha permesso, per la
prima volta, di confrontare direttamente il pavimento predetto con
l'errore misurato su un campione i.i.d.\ estratto dal modello esatto.

\paragraph{Documento di riferimento.}
\texttt{maxent\_pcd\_paper\_5\_arxiv.tex}, 2455 righe. I numeri di riga
citati in questa nota si riferiscono a quella versione.

\paragraph{Stato.} Il referee report di ACM TKDD non è ancora
pervenuto. Nessuna correzione è stata applicata al sorgente. Questa
nota registra la diagnosi e prepara il materiale di correzione, da
attivare quando si deciderà come procedere.

\paragraph{Co-autori.} La questione riguarda un lavoro co-firmato con
F.~Pachet e J.-D.~Zucker. Va comunicata prima di qualunque
risottomissione. \OPEN{predisporre versione inglese di questa nota}

% ======================================================================
\section{Definizione di MRE nel paper}

Il paper definisce (riga 791--792):
\begin{equation}
  \MRE_t \;=\; \frac{1}{m}\sum_{j=1}^{m}
  \frac{\bigl|\hat\alpha_j^{\,t}-\alpha_j\bigr|}{\alpha_j},
  \label{eq:mre}
\end{equation}
cioè la media degli errori \textbf{relativi in valore assoluto} sugli
$m$ vincoli di training. La definizione è esplicita e non ambigua; è
questa la quantità rispetto alla quale va confrontato qualunque
pavimento.

% ======================================================================
\section{L'errore}

\subsection{Cosa dice il paper}

Il pavimento è enunciato come
\begin{equation}
  \MRE_{\min} \;\approx\; \frac{1}{2\sqrt{N}} ,
  \label{eq:floor_paper}
\end{equation}
senza derivazione, nei seguenti punti:

\begin{center}
\begin{tabular}{ll}
\toprule
Riga & Contesto \\
\midrule
827        & Tabella \texttt{tab:guidelines}, colonna ``Rationale'' \\
1011--1012 & Didascalia figura: retta di riferimento $1/\sqrt{N}$ \\
1677       & ``adaptive rule before reaching the variance floor'' \\
1717--1721 & Didascalia \texttt{A\_ISTAT\_3\_pareto}: retta punteggiata \\
1858--1860 & Paragrafo \emph{Approximate convergence} \\
1872       & Stesso paragrafo, ``acceptable cost'' \\
\bottomrule
\end{tabular}
\end{center}

La formula $\alpha$-dipendente corretta \textbf{non compare in alcun
punto del paper}.

\subsection{Perché è sbagliata}

Sia $\hat\alpha_j$ la frazione empirica del vincolo $j$ su un campione
i.i.d.\ di taglia $N$, cosicché $N\hat\alpha_j\sim
\mathrm{Bin}(N,\alpha_j)$. L'errore standard \emph{assoluto} è
\begin{equation}
  \sigma_j \;=\; \sqrt{\frac{\alpha_j(1-\alpha_j)}{N}} ,
\end{equation}
il cui massimo su $\alpha_j$ si ha in $\alpha_j=1/2$ e vale
\begin{equation}
  \max_{\alpha}\sigma \;=\; \frac{1}{2\sqrt{N}} .
\end{equation}

L'espressione \eqref{eq:floor_paper} è dunque il \textbf{massimo errore
assoluto}. L'MRE di \eqref{eq:mre} divide invece per $\alpha_j$: si
tratta di un errore relativo, e le due quantità coincidono solo se
$\alpha_j\approx 1$. Per vincoli rari ($\alpha_j\ll 1/2$) la divisione
per $\alpha_j$ \emph{alza} il pavimento, anziché lasciarlo invariato.

Si tratta quindi di una confusione fra errore assoluto ed errore
relativo, non di un'approssimazione grossolana.

\subsection{Derivazione corretta}

Per $N\alpha_j\gtrsim 10$ vale l'approssimazione normale
$\hat\alpha_j\sim\mathcal{N}(\alpha_j,\sigma_j^2)$. Lo scarto assoluto
medio di una gaussiana è
\begin{equation}
  \mathbb{E}\bigl|\hat\alpha_j-\alpha_j\bigr|
  \;=\; \sqrt{\tfrac{2}{\pi}}\;\sigma_j ,
\end{equation}
da cui il contributo del singolo vincolo all'MRE
\begin{equation}
  \mathbb{E}\frac{|\hat\alpha_j-\alpha_j|}{\alpha_j}
  \;=\; \sqrt{\tfrac{2}{\pi}}\;
  \sqrt{\frac{1-\alpha_j}{\alpha_j N}} ,
\end{equation}
e infine il pavimento
\begin{empheq}[box=\fbox]{equation}
  \MRE_{\min}(N)
  \;=\; \sqrt{\tfrac{2}{\pi}}\;
  \frac{1}{m}\sum_{j=1}^{m}\sqrt{\frac{1-\alpha_j}{\alpha_j N}} .
  \label{eq:floor_ok}
\end{empheq}

\paragraph{Due osservazioni.}
\begin{enumerate}[leftmargin=*]
\item Il fattore $\sqrt{2/\pi}=0.797885$ è la conversione fra
  deviazione standard e scarto assoluto medio. È necessario perché
  \eqref{eq:mre} è definita con il valore assoluto e non in forma
  quadratica. Se in una versione precedente della formula lo si era
  omesso, la predizione risultava sovrastimata di un fattore $1.253$.
\item Il rapporto fra i due pavimenti è
  \begin{equation}
    \frac{\MRE_{\min}}{1/(2\sqrt{N})}
    \;=\; 2\sqrt{\tfrac{2}{\pi}}\;
    \frac{1}{m}\sum_j \sqrt{\frac{1-\alpha_j}{\alpha_j}}
    \;\approx\; 1.596\;\overline{\alpha^{-1/2}} ,
  \end{equation}
  indipendente da $N$. Dipende solo dalla distribuzione degli
  $\alpha_j$, cioè dalla rarità tipica dei vincoli.
\end{enumerate}

\paragraph{Condizione di validità.} L'approssimazione normale richiede
$N\alpha_j\gtrsim 10$. Con \texttt{--min-alpha 2e-4} e $N=10^5$ si ha
$N\alpha_{\min}\approx 21$: al limite ma accettabile. Per pool più
piccoli o soglie più basse la \eqref{eq:floor_ok} va usata con
cautela, o sostituita dallo scarto assoluto medio binomiale esatto.

% ======================================================================
\section{Validazione empirica}

Il test usa campioni i.i.d.\ di taglia $n$ estratti dal
\emph{modello esatto} (non dal pool PCD), su cui si misura l'MRE
rispetto ai target $\alpha_j$ e lo si confronta con
\eqref{eq:floor_ok}.

\begin{table}[h]
\centering
\caption{Pavimento predetto contro MRE misurato su campione i.i.d.\
dal modello esatto. Entrambe le popolazioni sono generate dalla
pipeline GSP; i due constraint set sono indipendenti per comune,
livello e cardinalità.}
\begin{tabular}{lrrrrrr}
\toprule
Dataset & $n$ & $m$ & $\MRE$ oss. & senza $\sqrt{2/\pi}$
        & con $\sqrt{2/\pi}$ & rapporto \\
\midrule
Parma K7C     & 198\,121 & 1\,094 & 0.0352 & 0.0448 & 0.0357 & \textbf{0.98} \\
Brescia K10C  & 198\,259 & 2\,611 & 0.0521 & 0.0660 & 0.0527 & \textbf{0.99} \\
\bottomrule
\end{tabular}
\label{tab:validation}
\end{table}

Il rapporto senza il fattore di conversione vale $0.79$ in entrambi i
casi, contro $\sqrt{2/\pi}=0.798$: l'accordo è esatto entro
l'arrotondamento. Con il fattore incluso l'accordo è all'1--2\% su due
dataset che differiscono per quattro ordini di grandezza in $|\mathcal{X}|$
($6.99\times10^4$ contro $3.73\times10^7$), per numero di attributi
($K{=}7$ contro $K{=}10$) e per numero di vincoli.

\medskip
\noindent
\textbf{Conclusione.} La \eqref{eq:floor_ok} è validata. La
\eqref{eq:floor_paper} sottostima il pavimento di un fattore $32$
(Parma) e $47$ (Brescia).

% ======================================================================
\section{Impatto sul paper}

\subsection{Cosa resta valido}

\begin{itemize}[leftmargin=*]
\item \textbf{Tutte le affermazioni sul tasso.} Entrambe le formule
  scalano come $O(1/\sqrt{N})$. Le rette di riferimento nei grafici
  log--log hanno la pendenza giusta; è sbagliata solo l'intercetta.
  Righe 915, 982, 1011, 1045: la sostanza tiene.
\item \textbf{La tesi centrale del paper.} Che $\MRE=0$ del raking sia
  algebricamente garantito e non una misura di accuratezza, e che
  diversità ($N_{\text{eff}}$, $H$) e valutazione held-out siano i
  criteri corretti, non dipende in alcun modo da questa formula.
\item \textbf{L'argomento del paragrafo \emph{Approximate
  convergence}} (1858--1872) ne esce \emph{rafforzato}: se il
  pavimento vero è $\sim\!5\times10^{-2}$ anziché $\sim\!10^{-3}$,
  l'inseguimento dell'MRE di training è ancora più chiaramente mal
  riposto.
\end{itemize}

\subsection{Cosa cade}

Il punto critico è la didascalia di riga 1717--1721:

\begin{quote}\itshape
the dotted line is the theoretical variance floor $1/(2\sqrt{N})$;
observed MRE lies above the floor at all $N$, confirming that mixing
bias dominates over gradient variance in this regime.
\end{quote}

Con un pavimento 30--50 volte troppo basso, \emph{qualunque} misura vi
sta sopra: il test non ha potere discriminante e l'inferenza sul
mixing bias non è sostenuta dall'evidenza addotta. Questo
indipendentemente dal fatto che la conclusione sia vera o falsa.

% ======================================================================
\section{La questione aperta: due regimi?}

\OPEN{determina la direzione della correzione}

Non è possibile stabilire il segno dell'effetto su A-ISTAT-3 senza
ricalcolare \eqref{eq:floor_ok} sui suoi $\alpha_j$. C'è però una
ragione strutturale per sospettare che A-ISTAT-3 e i casi censuari
stiano in regimi diversi:

\begin{table}[h]
\centering
\caption{Campioni di pool per vincolo. Il rapporto $N/m$ misura quanto
margine ha l'ottimizzatore per assecondare il pool su cui è calcolato
il gradiente.}
\begin{tabular}{lrrr}
\toprule
Esperimento & $m$ & $N$ & $N/m$ \\
\midrule
A-ISTAT-3 (paper) &    31 & 100\,000 & 3\,226 \\
Parma K7C         & 1\,094 & 100\,000 &     91 \\
Brescia K10C      & 2\,611 & 200\,000 &     77 \\
\bottomrule
\end{tabular}
\label{tab:regimi}
\end{table}

Con 31 gradi di libertà e $3\,200$ campioni per vincolo,
l'ottimizzatore ha pochissimo margine di sovradattamento: in
A-ISTAT-3 il mixing bias \emph{può} realmente dominare, e la
conclusione del paper essere corretta pur essendo sostenuta da un
argomento sbagliato. Nei casi censuari, con 40 volte meno campioni per
vincolo, il sovradattamento in-sample è plausibile.

Se la lettura a due regimi è confermata, la correzione non è una
ritirata ma un risultato aggiuntivo, con $N/m$ come variabile di
controllo.

% ======================================================================
\section{Evidenza collaterale: l'MRE del Gibbs non è errore campionario}

Confrontando l'MRE del \emph{fit} PCD con il pavimento corretto
calcolato alla taglia del pool:

\begin{table}[h]
\centering
\caption{L'MRE del Gibbs sta \emph{sotto} il pavimento di
campionamento. La taglia i.i.d.\ equivalente
$N_{\text{eff}}=N(\MRE_{\min}/\MRE_{\text{oss}})^2$ eccede $N$, il che
è impossibile per un campionamento onesto e a maggior ragione per una
catena autocorrelata.}
\begin{tabular}{lrrrrr}
\toprule
Dataset & $N$ pool & $\MRE$ Gibbs & $\MRE_{\min}(N)$ & rapporto
        & $N_{\text{eff}}$ apparente \\
\midrule
Parma K7C    & 100\,000 & 0.0396 & 0.0503 & 0.79 & $\approx161\,000$ \\
Brescia K10C & 200\,000 & 0.0323 & 0.0524 & 0.62 & $\approx526\,000$ \\
\bottomrule
\end{tabular}
\end{table}

\paragraph{Interpretazione.} Il pool PCD non è un campione i.i.d.\
estratto dal modello: è l'oggetto \emph{rispetto al quale} i $\lambda$
vengono ottimizzati. Il gradiente è (target $-$ momento del pool), e le
iterazioni esterne spingono proprio quella differenza verso zero. L'MRE
riportato è dunque un \textbf{residuo di ottimizzazione misurato sullo
stesso campione usato per ottimizzare}, non un errore di
campionamento. Non ha un pavimento e, in linea di principio, tende a
zero con le iterazioni.

\paragraph{Conferma indipendente.} Le divergenze contro il modello
esatto restano grandi e asimmetriche mentre l'MRE è basso:

\begin{center}
\begin{tabular}{lrrr}
\toprule
Dataset & $\MRE$ esatto & $\mathrm{KL}(\text{exact}\|\text{gibbs})$
        & $\mathrm{KL}(\text{gibbs}\|\text{exact})$ \\
\midrule
Parma K7C    & $4.19\times10^{-4}$ & 0.1056 & 0.6415 \\
Brescia K10C & $4.61\times10^{-4}$ & 0.1041 & 0.7406 \\
\bottomrule
\end{tabular}
\end{center}

L'asimmetria $\approx6$--$7\times$ è stabile fra i due dataset. Poiché
$\mathrm{KL}(g\|e)=\sum g\log(g/e)$ esplode dove $e\to0$ e $g$ no, il
pool PCD sta mettendo massa dove il modello esatto non ne mette: è
\emph{sovradispersione}, non sottocampionamento. Coerente con i
supporti misurati su Parma (esatto $52\,167/69\,888$, Gibbs
$67\,704/69\,888$).

\paragraph{Nota metodologica.} È la stessa critica che il paper già
muove al raking --- l'MRE di training non misura accuratezza ---
applicata però al metodo proprio. Presentata così, rafforza la
credibilità dell'argomento in revisione anziché indebolirla.

% ======================================================================
\section{Verifiche da eseguire}

\begin{enumerate}[leftmargin=*]
\item \textbf{Ricalcolo del pavimento sugli esperimenti del paper.}
  Applicare \eqref{eq:floor_ok} agli $\alpha_j$ di Syn-ISTAT
  (A-ISTAT-1/2/3) e confrontare con gli MRE già in tabella. Determina
  il segno della correzione. \emph{Prerequisito: i constraint set di
  Syn-ISTAT devono essere ancora su disco.}
\item \textbf{Pendenza delle curve di convergenza.} Verificare su
  \texttt{gibbs\_mre\_curve} se i fit confrontati sono a plateau o
  ancora in discesa. Pendenza log--log $\approx-0.5$ indica
  convergenza in corso, $\approx0$ plateau. Fit troncati a stadi
  diversi non sono confrontabili.
\item \textbf{Traiettoria congiunta MRE--KL.} Salvare
  $\mathrm{KL}(\text{gibbs}\|\text{exact})$ ogni 50 iterazioni
  esterne. \emph{MRE che scende mentre la KL sale} sarebbe la
  dimostrazione diretta del sovradattamento in-sample, e diventerebbe
  un risultato centrale anziché una diagnostica.
\item \textbf{Sensibilità a \texttt{--min-alpha}.} Il 10\% di vincoli
  più rari contribuisce il 25\% (Parma) e il 21\% (Brescia) del
  pavimento. Fit con soglie diverse hanno pavimenti diversi e non sono
  confrontabili: verificare che nel paper le soglie siano uniformi
  fra le righe delle tabelle comparative.
\end{enumerate}

% ======================================================================
\section{Materiale di correzione (inglese, pronto all'uso)}

\selectlanguage{english}
\noindent
\emph{Da attivare solo dopo la verifica 1. Il testo per la didascalia
di riga 1717 è dato in due varianti alternative.}

\subsection*{Definizione del pavimento (nuovo, da inserire in Sez.~4)}

\begin{lstlisting}[language=TeX]
For an i.i.d. sample of size $N$ from the fitted model,
$N\hat\alpha_j \sim \mathrm{Bin}(N, \alpha_j)$, so the absolute
standard error of $\hat\alpha_j$ is
$\sigma_j = \sqrt{\alpha_j(1-\alpha_j)/N}$.
Since $\MRE$ in \cref{eq:mre} is a mean of \emph{relative}
absolute deviations, and $\mathbb{E}|\hat\alpha_j - \alpha_j|
= \sqrt{2/\pi}\,\sigma_j$ under the normal approximation
($N\alpha_j \gtrsim 10$), the sampling floor is
\begin{equation}
  \MRE_{\min}(N) = \sqrt{\frac{2}{\pi}}\,
  \frac{1}{m}\sum_{j=1}^{m}\sqrt{\frac{1-\alpha_j}{\alpha_j N}}.
  \label{eq:mre_floor}
\end{equation}
Note that $\MRE_{\min}$ depends on the rarity of the constraints
and not on $N$ alone: for census marginals with
$\alpha_j \ll 1/2$ it exceeds the naive bound $1/(2\sqrt{N})$
--- the maximum \emph{absolute} standard error --- by one to two
orders of magnitude. We verify \cref{eq:mre_floor} on two
independent synthetic populations, obtaining agreement within
$2\%$ (\cref{tab:floor_validation}).
\end{lstlisting}

\subsection*{Didascalia riga 1717 --- variante A (se l'osservato resta sopra)}

\begin{lstlisting}[language=TeX]
the dotted line is the sampling floor $\MRE_{\min}(N)$ of
\cref{eq:mre_floor}, evaluated on the constraint targets;
observed MRE lies above the floor at all $N$, confirming that
mixing bias dominates over gradient variance in this regime.
\end{lstlisting}

\subsection*{Didascalia riga 1717 --- variante B (se l'osservato scende sotto)}

\begin{lstlisting}[language=TeX]
the dotted line is the sampling floor $\MRE_{\min}(N)$ of
\cref{eq:mre_floor}. Observed training MRE falls below this
floor, which is expected: the pool is the sample against which
$\blam$ is optimised, so training MRE is an in-sample residual
rather than a sampling error. It is therefore not a valid
convergence criterion, and we report KL divergence against the
exact solution instead (\cref{tab:kl}).
\end{lstlisting}

\subsection*{Tabella guidelines, riga 827}

\begin{lstlisting}[language=TeX]
  & Sampling floor $\MRE_{\min}(N)$ grows with constraint
    rarity; sweet spot at $N{=}25$K \\
\end{lstlisting}

\subsection*{Paragrafo \emph{Approximate convergence}, riga 1858}

\begin{lstlisting}[language=TeX]
GibbsPCDSolver does not guarantee exact satisfaction of the
constraints. The relevant floor is $\MRE_{\min}(N)$ of
\cref{eq:mre_floor}, which for census-scale constraint sets
is of order $5\times10^{-2}$ at $N{=}10^5$ --- one to two
orders of magnitude above the naive $1/(2\sqrt{N})$ bound.
\end{lstlisting}

\selectlanguage{italian}

% ======================================================================
\section{Codice di riproduzione}

\begin{lstlisting}[language=Python]
import json, os
import numpy as np, pandas as pd

SQRT_2_PI = np.sqrt(2 / np.pi)   # 0.797885

def mre_floor(comune, livello, anno=2024):
    """Pavimento predetto vs MRE su campione i.i.d. dal modello esatto."""
    d = os.path.expanduser(
        f"~/progetti/gsp/data/comuni/{comune}/constraints_{anno}")
    cs  = json.load(open(f"{d}/cs_{livello}.json"))
    fit = json.load(open(f"{d}/fit_{livello}.json"))
    pop = pd.read_csv(f"{d}/popolazione_{livello}.csv", dtype=str)

    V, C = cs["vars"], cs["categories"]
    code = {v: pop[v].astype(str)
                 .map({c: i for i, c in enumerate(C[v])}).to_numpy()
            for v in V}

    amap = {(tuple(c["attrs"]), tuple(c["vals"])): c["alpha"]
            for c in cs["constraints"]}
    kept = [(tuple(a), tuple(b)) for a, b in fit["kept_sig_exact"]]
    n = len(pop)

    a, emp = [], []
    for at, va in kept:
        m = np.ones(n, dtype=bool)
        for j, val in zip(at, va):
            m &= (code[V[j]] == val)
        a.append(amap[(at, va)]); emp.append(m.mean())
    a, emp = np.array(a), np.array(emp)

    mre   = float(np.mean(np.abs(emp - a) / a))
    floor = SQRT_2_PI * float(np.mean(np.sqrt((1 - a) / (a * n))))
    print(f"{comune} {livello}: n={n:,} m={len(a)} "
          f"MRE={mre:.4f} floor={floor:.4f} ratio={mre/floor:.2f}")

    # pavimento alla taglia del pool, per confronto con l'MRE del Gibbs
    N = fit["pool"]
    fN = SQRT_2_PI * float(np.mean(np.sqrt((1 - a) / (a * N))))
    g  = fit["gibbs"]["mre_pos"]
    print(f"   pool N={N:,}: MRE_gibbs={g:.4f} floor={fN:.4f} "
          f"ratio={g/fN:.2f} -> N_eff apparente {N*(fN/g)**2:,.0f}")

mre_floor("034027", "K7C")
mre_floor("017029", "K10C")
\end{lstlisting}

% ======================================================================
\section*{Changelog}
\addcontentsline{toc}{section}{Changelog}

\begin{description}[leftmargin=2.2cm,style=nextline]
\item[v0.1 --- 27/07/2026] Prima stesura. Diagnosi completa,
  validazione su Parma K7C e Brescia K10C, materiale di correzione
  predisposto. Verifiche 1--4 non ancora eseguite. Nessuna modifica
  applicata al sorgente del paper.
\end{description}

\end{document}
